Большинство архивных химических книг и сканированных материалов характеризуются неоднородным качеством. В работе~\cite{refAkhil} выделялись следующие проблемы качества сканированных книг:
\begin{itemize}
	\item Низкое разрешение / размытость
	\item Шум, зернистость, артефакты сжатия
	\item Неравномерная засветка / тени
	\item Ориентация страницы с малым поворотом
\end{itemize}

Эти факторы напрямую влияют на точность распознания как и графических формул и реакций так и текста. Поэтому перед обработкой самой OCR необходимо сделать пред обработку и нормализацию изображения. 



